Newly acquired memories are not fixed at the moment of encoding. They remain
susceptible to interference and undergo post-learning processes that stabilize,
reorganize, and integrate them into longer-term representations. Sleep supports
this transformation through coordinated neural activity during non-rapid eye
movement (NREM) sleep, including the temporal interaction of slow oscillations,
sleep spindles, and hippocampal sharp-wave ripples
\citep{diekelmann2010memory,rasch2013sleep,klinzing2019mechanisms}.
These processes contribute to the reactivation and reorganization of recently
encoded information rather than merely protecting memory from continued waking
interference.

Targeted Memory Reactivation (TMR) provides an experimental means of influencing
this endogenous consolidation process. In a typical TMR paradigm, sensory cues
are associated with specific learning material during wakefulness and are
subsequently re-presented during sleep to bias reactivation toward the
corresponding memories \citep{oudiette2013upgrading}. Across experimental
paradigms, appropriately delivered cues have produced a modest average benefit
for later memory, with a meta-analytic effect size of Hedges'
\(g=0.29\). The magnitude and reliability of this effect nevertheless vary with
memory domain, encoding quality, cueing protocol, and physiological context
\citep{hu2020promoting}. TMR should therefore be understood as a
condition-dependent modulation of memory consolidation rather than as a general
or guaranteed method of memory enhancement.

The principal translational difficulty is that experimental efficacy does not
by itself define a practical intervention system. Laboratory TMR commonly
depends on physiological recording, controlled cue timing, continuous
observation, and the ability to suspend stimulation when sleep becomes unstable.
A system intended to operate outside this environment must reproduce the
intervention-relevant functions of that workflow rather than merely generate a
retrospective description of the night. It must acquire physiological
information during ongoing sleep, determine whether the current interval is
suitable for cueing, execute stimulation with bounded and measurable timing,
and evaluate whether the cue has disrupted the state it was intended to exploit.

Recent studies provide proof-of-concept evidence that portions of this process
can be automated in home environments. SleepStim used smartwatch-derived
movement and cardiac information to guide auditory cue delivery, demonstrating
the feasibility of peripheral-sensor-based automated TMR while also showing that
outcomes depend on stimulation conditions and the avoidance of sleep disruption
\citep{whitmore2022improving}. A subsequent study used a wearable EEG
headband to deliver vocabulary cues during detected slow-wave activity and
reported a benefit for cued material under home conditions
\citep{salfi2025promoting}. These studies narrow the distance between
laboratory TMR and wearable implementation, but they do not establish that all
sensing approaches provide equivalent control reliability or behavioral
efficacy.

A critical distinction therefore remains between general sleep assessment and
intervention-ready physiological monitoring. A sensing modality may classify
sleep adequately in retrospective analysis while lacking the continuous signal
access, causal state estimation, uncertainty handling, cue synchronization, or
post-stimulation monitoring required for real-time closed-loop operation.
Wearable EEG and peripheral sensing, including photoplethysmography (PPG), should
consequently be evaluated not only by their agreement with conventional sleep
staging, but by whether they can support the physiological decisions on which
effective and interpretable TMR depends. The relevant translational problem is
not whether the necessary components exist independently, but whether the
available evidence supports their combination within a coherent
first-generation pathway without converting unresolved assumptions into system
requirements.

This paper addresses the following central research question:

\begin{quote}
\textit{How can current evidence on sleep-dependent memory consolidation,
Targeted Memory Reactivation, and wearable sleep monitoring be translated into
a scientifically defensible first-generation pathway for real-time closed-loop
TMR?}
\end{quote}

To answer this question, the study develops and applies a structured
evidence-to-engineering methodology. Biological and experimental findings are
first translated into system-level requirements, after which candidate
EEG-based and peripheral-sensor-based implementation pathways are evaluated
according to physiological relevance, temporal suitability, evidence maturity,
and validation risk. The intended contribution is not a new biological
mechanism, engineering algorithm, or validated product. It is a traceable
translational analysis that distinguishes established findings from unresolved
implementation assumptions and identifies the first-generation direction most
strongly justified for subsequent investigation.

The remainder of the paper is organized as follows. Section~2 defines the
research scope and evidence-to-engineering methodology. Section~3 introduces the
neuroscientific mechanisms required to interpret sleep-dependent memory
processing, and Section~4 examines the biological rationale and experimental
evidence for TMR. Section~5 evaluates wearable sleep monitoring in relation to
real-time intervention requirements. Section~6 translates the combined evidence
into system-level requirements and compares first-generation implementation
pathways. Sections~7 and~8 define the limitations and future research
directions, respectively, and Section~9 presents the final conclusions.